\chapter{Модуль и числовые промежутки}
\begin{definition}\term{Модуль} (\emph{absolute value}) \(| x |\) — расстояние от числа \(x\) до нуля на числовой прямой: \(| x | = x\) при \(x \ge  0\) и \(| x | = - x\) при \(x < 0\).\end{definition}
\begin{definition}\term{Числовой промежуток} (\emph{interval}) — множество всех чисел между заданными границами с указанным включением или исключением концов.\end{definition}
\begin{definition}\term{Пересечение} (\emph{intersection}) множеств \(A\) и \(B\), обозначаемое \(A\cap B\), — множество элементов, принадлежащих одновременно обоим множествам.\end{definition}
\begin{definition}\term{Объединение} (\emph{union}) множеств \(A\) и \(B\), обозначаемое \(A\cup B\), — множество элементов, принадлежащих хотя бы одному из этих множеств.\end{definition}
\section*{Свойства модуля}
\(| x | \ge  0;\allowbreak | x | = 0\) только при \(x = 0\). Также \(| - x | = | x |\) и \(| ab | = | a | | b |\).\par
Расстояние между числами \(a\) и \(b\) равно \(| a - b |\).\par
Неравенство треугольника: \(| a + b | \le  | a | + | b |\).\par
При \(r > 0\) условие \(| x - a | < r\) означает \(a - r < x < a + r\). При \(r>0\) условие \(|x-a|\le r\) равносильно \(a-r\le x\le a+r\).\par
\section*{Запись промежутков}
\((a;\allowbreak b)\): \(a < x < b\). \([a;\allowbreak b]\): \(a \le x \le b\). \([a;\allowbreak b)\) и \((a;\allowbreak b]\) включают только один из концов.\par
Бесконечные границы всегда записывают с круглой скобкой: например, \((- \infty ;\allowbreak 3]\).\par
{\small\textit{Пример: \(| x - 2 | \le  3\) равносильно \(- 1 \le  x \le  5\), то есть \(x \in  [- 1;\allowbreak 5]\).}}\par

\Needspace{7\baselineskip}\section*{Задачи}
\begin{enumerate}[label={},leftmargin=3.2em,itemsep=10pt]
\Needspace{3\baselineskip}\item[\textbf{8.1.}] Вычисли \(| - 8 | + | 3 |\).
\Needspace{3\baselineskip}\item[\textbf{8.2.}] Вычисли \(| - 5 + 2 |\).
\Needspace{3\baselineskip}\item[\textbf{8.3.}] Найди расстояние между \(- 7\) и \(4\).
\Needspace{3\baselineskip}\item[\textbf{8.4.}] Сравни \(| - 6 |\) и \(| 5 |\).
\Needspace{3\baselineskip}\item[\textbf{8.5.}] Запиши \(- 2 < x \le  3\) промежутком.
\Needspace{3\baselineskip}\item[\textbf{8.6.}] Запиши \([- 4;\allowbreak 1)\) двойным неравенством.
\Needspace{3\baselineskip}\item[\textbf{8.7.}] Выпиши целые числа из \((- 3;\allowbreak 2]\).
\Needspace{3\baselineskip}\item[\textbf{8.8.}] Найди пересечение \([- 2;\allowbreak 4]\) и \((1;\allowbreak 7)\).
\Needspace{3\baselineskip}\item[\textbf{8.9.}] Найди объединение \((- \infty ;\allowbreak 2]\) и \((1;\allowbreak 5)\).
\Needspace{3\baselineskip}\item[\textbf{8.10.}] Реши \(| x | = 5\).
\Needspace{3\baselineskip}\item[\textbf{8.11.}] Реши \(| x - 2 | < 3\).
\Needspace{3\baselineskip}\item[\textbf{8.12.}] Упрости \(\sqrt{x^{2}}\) и вычисли при \(x = - 4\).
\Needspace{3\baselineskip}\item[\textbf{8.13.}] Из жизни. Температура в морозилке \(- 18 {}^{\circ} C\), а в комнате \(+ 22 {}^{\circ} C\). На сколько градусов отличаются температуры?
\Needspace{3\baselineskip}\item[\textbf{8.14.}] Из жизни. Длина детали должна отличаться от \(50\) мм не более чем на \(0,2\) мм. Запиши допустимые длины \(x\) через модуль и промежутком.
\Needspace{3\baselineskip}\item[\textbf{8.15.}] Из жизни. Аттракцион доступен при росте от \(120\) см включительно до \(190\) см не включительно. Запиши условие для роста \(h\) неравенством и промежутком. Можно ли войти при росте \(190\) см?
\Needspace{3\baselineskip}\item[\textbf{8.16.}\starmark] Реши \(| x - 3 | + | x + 1 | = 6\).
\Needspace{3\baselineskip}\item[\textbf{8.17.}\starmark] Найди все целые \(x\), для которых \(| 2 x - 1 | \le  5\) и \(x \notin  [0;\allowbreak 2)\).
\Needspace{3\baselineskip}\item[\textbf{8.18.}\starmark] Два датчика допускают температуру: первый \(| t - 20 | \le  3\), второй \(| t - 24 | < 4\). Найди общий допустимый диапазон.
\Needspace{3\baselineskip}\item[\textbf{8.19.}] Докажите обратное неравенство треугольника \(\bigl||a|-|b|\bigr|\leq|a-b|\) для действительных \(a,b\).
\Needspace{3\baselineskip}\item[\textbf{8.20.}] Докажите, что при \(r>0\) неравенство \(|x-a|\leq r\) равносильно \(a-r\leq x\leq a+r\).
\end{enumerate}
